\section{Worum es geht}

\textbf{[SETZUNG]} Die Zeit ist in FFGFT keine äußere Bühne, sondern Teil der
Geometrie. Aus der Zeit-Masse-Dualität $\tilde T\cdot m = 1$ (A020) folgt in
natürlichen Einheiten die native Zeit-Energie-Reziprozität $T\cdot E = 1$ -- eine
\emph{Gleichung}, keine Ungleichung. Dieses Dokument entwickelt daraus vier
Dinge: die \emph{Verortung} der Zeit (im Zustandsraum statt daneben), die
\emph{Reziprozität} und ihre Ränder (kleinstes Zeitinkrement, Schwarzes Loch als
Extremum, kein leerer Zustand), den \emph{Grundtakt} (die zulässigen Moden sind
kommensurable Vielfache eines einzigen $\xi$-fixierten Inkrements) und den
\emph{Zeitpfeil} (geometrisch, nicht thermodynamisch primär). Der kosmisch-große
Sektor (Skalen jenseits der Teilchenphysik) wird hier nicht geführt (A010); er
verlangt eine eigene, unabhängig gesetzte Skala -- warum, steht unten.

\section{Die Gabelung: Zeit im Zustandsraum oder daneben}

\textbf{[B]} Am Anfang jeder quantenmechanischen Modellbildung steht eine meist
unausgesprochene Wahl: \emph{Ist die Zeit eine Richtung im Zustandsraum, oder ein
Parameter neben ihm?} Die Frage ist nicht technisch, sondern ontologisch. Beide
Antworten arbeiten mit demselben Formalismus -- selbstadjungierte Operatoren,
Spektren, Projektoren --; der Unterschied liegt allein darin, \emph{wo} die Zeit
sitzt. Und genau das entscheidet, ob die physikalischen Größen aus der Struktur
\emph{fallen} oder ihr \emph{angeheftet} werden müssen.

\textbf{[B]} Die Frage reicht bis zu den Wurzeln der Quantentheorie. Schon 1932
erkannte Pauli, dass ein Zeitoperator, der kanonisch zu einem nach unten
beschränkten Hamiltonoperator konjugiert wäre, nicht existiert (Pauli-Theorem);
das machte die Zeit als äußeren Parameter scheinbar unumstößlich. Spätere Arbeiten
(Aharonov--Bohm, Garrison--Wong) zeigten, dass unter bestimmten Bedingungen --
etwa bei nicht halb-beschränkten Spektren -- Zeitoperatoren doch konstruierbar
sind. Der Streit ist bis heute nicht beigelegt und bildet den Hintergrund der
Alternative.

\section{Die ausgerollte Sicht und ihr Preis}

\textbf{[B]} Die vertraute Formulierung siedelt die Zeit außerhalb des
Zustandsraums an: $\mathcal H = L^2(\text{Konfigurationsraum})$ ist der Raum der
Zustände \emph{zu einem Zeitpunkt}, die Zeit der Parameter $t$ der unitären
Evolution $e^{-i\hat H t}$. Observablen -- Ort, Impuls, Spin -- sind Operatoren
\emph{auf} $\mathcal H$; die Zeit ist keiner (Pauli). Die Folge ist strukturell:
weil die Zeit außen steht, enthält der Zustandsraum allein keine dynamischen
Skalen. Massen, Kopplungen, Lebensdauern sind in der reinen Raumstruktur nicht
vorhanden -- sie werden über zusätzliche Strukturen eingeführt, im Standardmodell
durch Yukawa-Terme mit freien Parametern. Die Struktur liefert den Rahmen, die
Zahlen kommen aus der Messung.

\textbf{[S]} Diese Kritik ist alt: schon Weyl und Wigner sahen die Massen
gruppentheoretisch als Casimir-Invarianten -- strukturelle Label, keine
angehefteten Parameter --, doch die Herkunft der Massenskalen blieb offen. Der
Punkt ist nicht, dass die ausgerollte Sicht falsch wäre, sondern dass sie diese
Frage unbeantwortet lässt.

\textbf{[S]} \textbf{Warum das Anheften ein Preis ist.} Viele Programme mit
geometrischem Anspruch folgen einem gemeinsamen Schema: (1) ein diskreter Träger
(Gitter, Cut-and-Project-Struktur, periodischer Graph); (2) eine
selbstadjungierte Operatorfamilie darauf; (3) Spektralprojektoren extrahieren
Sektoren; (4) diesen Sektoren wird \emph{nachträglich} Bedeutung zugeordnet --
Masse häufig über eine parametrische Abbildung mit ein oder zwei freien
Parametern. Ein solches Vorgehen kann mathematisch vollständig kontrolliert sein;
entscheidend ist seine \emph{Verortung der Zeit}: der Träger ist räumlich, die
Sektoren sind räumliche Eigenräume, keine Zeit-Moden -- die Masse ist nicht
enthalten, sondern wird angeheftet. Und hier greift ein Argument, das sich selbst
trägt: eine parametrische Abbildung mit freien Parametern kann jede endliche
Werteliste treffen. Ist sie nicht vorab und universell fixiert, benennt sie die
bekannten Eingaben nur um. Um Zahlen zu liefern, muss sie an externe Werte
gekoppelt werden -- an die Messung, und damit an das Standardmodell. Ein Beispiel
aus der etablierten Physik: die Graphen-artige Modellierung definiert
Quasiteilchen auf hexagonalen Gittern mit linearer Dispersion (masselose
Fermionen); eine Masse erfordert stets zusätzliche Terme (Haldane, Kekulé) mit
freien Parametern. Die Geometrie liefert die Kinematik, nicht die Massenskala. Die
\emph{Schwierigkeit} dieses Wegs ist kein Zufall, sondern die unmittelbare Folge
der ausgerollten Verortung der Zeit.

\section{Die eingerollte Sicht: Zeit als kompakte Dimension}

\textbf{[B]} Die Alternative belässt die Zeit nicht außen und eliminiert sie auch
nicht, sondern zieht sie \emph{in} den Zustandsraum -- als geschlossene,
periodische Dimension. An die Stelle von $L^2(\text{Raum})$ mit äußerem $t$ tritt
ein Raum, in dem eine kompakte Zeit-Richtung Teil der Geometrie ist; im Korpus ist
das $\mathcal H = L^2(T^4)\otimes\mathbb{C}^3$ (A070), eine Torus-Richtung trägt
die Zeit. Auf einer geschlossenen Dimension sind die zulässigen Zustände
\emph{stehende Moden mit diskreten Windungszahlen}: die Quantisierung ist nicht
postuliert, sondern geometrisch erzwungen. Die Masse wird zum \emph{Eigenwert
eines Operators auf $\mathcal H$ selbst} -- nicht angeheftet, sondern strukturell
enthalten. Die Verhältnisse der Massen liegen fest, weil die zulässigen Moden
festliegen.

\textbf{[B]} Konkret entwickelt man eine Funktion auf einer kompakten
Zeit-Richtung der Länge $T$ in eine Fourier-Reihe,
\[
  \Psi(t,\mathbf x) = \sum_{n\in\mathbb Z} e^{i n t/T}\,\psi_n(\mathbf x),
\]
die diskreten Frequenzen $\omega_n = 2\pi n/T$ sind (in natürlichen Einheiten) die
Massenwerte, erzeugt vom Operator $\hat E = -i\,\partial_t$ mit Spektrum
$\{2\pi n/T\}$. Die Massenquantisierung ist eine direkte Folge der Periodizität --
nichts muss postuliert werden. Die Kopplungen zwischen Moden,
\[
  g_{ijk} \propto \int e^{i(n_i-n_j-n_k)t/T}\,dt = T\,\delta_{n_i,\,n_j+n_k},
\]
sind exakte Auswahlregeln für die Moden-Indizes -- die \enquote{Yukawa-Kopplungen}
erscheinen als Erhaltungsregeln, nicht als freie Zahlen.

\textbf{[B]} Dass eine kompakte Dimension einen diskreten Moden-Turm mit
$m_n\propto n/R$ erzeugt, ist der Mechanismus von Kaluza und Klein -- hier für die
eingerollte \emph{Zeit} statt für eine Raum-Richtung, mit der Masse in der Rolle
des Modenindex. Und Paulis Einwand trifft diese Sicht \emph{nicht}: er gilt für
eine offene, nach unten unbeschränkte Zeit; eine kompakte Zeit verhält sich wie
eine Winkelvariable, ihr konjugiertes Spektrum ist diskret (wie Drehimpuls zum
Winkel), das Pauli-Argument entfällt. Die geschlossene Zeit ist der natürliche Ort
einer diskreten Modenstruktur.

\section{Zeit-Masse-Dualität und native Reziprozität}

\textbf{[B]} Die konkrete Ausführung ist $\tilde T\cdot m = 1$: jede massive Größe
\emph{ist} ihr eigener Zeitzyklus. In natürlichen Einheiten ($E=m$) folgt
\[
  \boxed{\;T\cdot E = 1\;}
\]
Das ist eine \emph{Gleichung}, keine Ungleichung -- eine schärfere Aussage als
Heisenbergs $\Delta E\cdot\Delta t\ge 1/2$, und eine korpuseigene. Sie ersetzt
eine früher geliehene Begründung: die Abwesenheit eines singulären Anfangs war im
Korpus einmal mit Heisenbergs Energie-Zeit-Unschärfe begründet. Diese Begründung
wird \emph{ersetzt}, nicht widerlegt -- die Konklusion bleibt, die Ableitung wird
nativ. Drei Mängel der geliehenen Fassung:

\begin{itemize}[leftmargin=2em,itemsep=3pt,topsep=3pt]
  \item \textbf{Invertierte Implikation.} Aus $\Delta E\ge 1/(2\Delta t)$ folgt
    bei \emph{endlichem} $\Delta t$ eine \emph{endliche} Schranke, nicht
    $\Delta E\to\infty$; letzteres folgt allein aus $\Delta t\to0$, also aus dem
    Augenblick $t=0$, nicht aus einem endlichen Anfang.
  \item \textbf{Anwendungsfehler.} Nach Paulis Theorem gibt es keinen
    selbstadjungierten Zeitoperator zu einem beschränkten Hamiltonoperator; die
    strenge Form (Mandelstam--Tamm) definiert $\Delta t$ als charakteristische
    \emph{Änderungszeit einer Observablen}, nicht als verstrichene Dauer. Ein
    \enquote{Alter} ist kein zulässiges $\Delta t$ dieser Relation.
  \item \textbf{Stilbruch.} \enquote{Beweist unwiderlegbar} ist eine
    Behauptungsstärke, die der Anspruchsledger nicht kennt; er kennt \emph{aus
    $\xi$ abgeleitet}, \emph{eingeschränkt}, \emph{blockiert}.
\end{itemize}

\textbf{[B]} Der übergreifende Punkt: Heisenbergs Relation ist im QM-Rahmen nicht
falsch, aber \emph{projiziert} -- sie gehört zur Beschreibungsebene, die FFGFT als
Projektion aus der tieferen Struktur versteht (A090). Ein Argument, das sich auf
sie stützt, verlässt die Projektion nicht, sondern zirkuliert in ihr.

\section{Das kleinste Inkrement}

\textbf{[K]} FFGFTs Zeit ist der kumulierte Defekt der ausgerollten Windung,
\[
  D(k) = \sum_{n\le k} 100\,\xi_n \;\to\; \ln\!\Bigl(1+\tfrac{k}{75}\Bigr),
  \qquad \xi_{n+1} = \xi_n\,(1-100\,\xi_n),
\]
mit dem ersten Schritt
\[
  d_1 = 100\,\xi_1 = 100\cdot\tfrac{4}{30000} = \tfrac{1}{75}.
\]
$D(k)$ wächst für $k\to\infty$ unbeschränkt (logarithmisch); die entscheidende
Aussage ist nicht \enquote{$D$ divergiert nie}, sondern: $D(k)$ ist an jeder
endlichen Stelle endlich, und kein Zeitinkrement wird kleiner als $d_1$. Aus
$T\cdot E=1$ folgt $E\to\infty\Leftrightarrow T\to0$; ein Zustand mit $T=0$
existiert nicht, weil die Zeit nach unten durch das endliche, $\xi$-fixierte
Inkrement beschränkt ist. Damit ist $E$ nach oben beschränkt -- eine unendliche
Energiedichte ist \emph{strukturell} ausgeschlossen, nicht durch eine Aussage über
Messgenauigkeit.

\section{Beide Ränder: Schwarzes Loch als Extremum, dünnster Zustand}

\textbf{[K]} Dieselbe Kette bestimmt das Schwarze Loch -- und das Ergebnis ist
unspektakulär: nichts divergiert. Aus $T\cdot E=1$ mit $T\ge \tau_0$ folgt
$E\le 1/\tau_0$. Die Extrema sind aus $\xi$ beziffert: mit $L_0=\xi\,\ell_P$ (A210)
liefert die Energie-Längen-Reziprozität
\[
  E_\text{max} = \frac{\hbar c}{L_0} = \frac{E_P}{\xi},
  \qquad \tau_0 = \xi\,t_P\quad(\text{minimale Zeitskala; per }c{=}1\text{ dieselbe Größe wie }L_0\text{ in Zeit-Kleidung, A080}),
\]
und die Reziprozität schließt \emph{exakt}, weil $\xi$ sich herauskürzt:
\[
  \tau_0\cdot E_\text{max} = (\xi\,t_P)\cdot\frac{E_P}{\xi} = t_P\cdot E_P = 1.
\]
Derselbe Cutoff erscheint in der Feldtheorie als natürlicher UV-Cutoff bei
$E_P/\xi$: die Geometrie stellt den Cutoff, ein nachträglicher Regulator ist nicht
nötig (A190).

\textbf{[K]} \textbf{Skalen, nicht Summen.} $\tau_0$ und $E_\text{max}$ sind Skalen,
keine Schranken für Gesamtheiten: $\tau_0$ ist das minimale Zeit\emph{inkrement} der
Windung, $E_\text{max}$ die maximale Energie \emph{pro Anregung} (ein Cutoff, der
stets eine einzelne Mode begrenzt, nie eine Summe). Die Gesamtenergie eines
zusammengesetzten Systems darf beliebig groß sein -- es besteht aus vielen Zellen,
nicht aus einer dichteren; \enquote{das $T$ der Erde} ist keine wohlgestellte
Größe. Der Horizont ist entsprechend der Ort, an dem die \emph{lokale} Dichte den
Zellen-Cutoff erreicht: langsamste Zeit $T=\tau_0$ (nicht $T\to0$), maximale
Konzentration $E=1/\tau_0$ pro Zelle, maximale Krümmung -- alles endlich und
wohldefiniert. Er ist kein Vorhang vor einer Singularität, sondern die Fläche, auf
der Zeit ihr Minimum und Konzentration ihr Maximum annimmt. In FFGFT gibt es
darum \emph{nichts zu reparieren}: Modelle mit Bounce oder Übergangsschale beheben
einen Defekt, den FFGFT gar nicht erzeugt.

\textbf{[S]} \textbf{Das andere Ende.} $T\cdot m=1$ hat bei $m=0$ keine Lösung --
FFGFT kennt keinen \emph{leeren} Zustand, nur einen dünnsten. Der scheinbare
Zirkel \enquote{Vakuum ist leer, wirkt nur durch Krümmung, Krümmung setzt Masse
voraus} entsteht aus einem kausalen Pfeil der Allgemeinen Relativität ($T_{\mu\nu}$
\emph{erzeugt} Krümmung). FFGFT hat diesen Pfeil nicht: $T\cdot m=1$ ist keine
Ursache-Wirkung-Beziehung, sondern eine Reziprozität -- Masse und zeitliche
Krümmung sind zwei Lesarten derselben Struktur, keine erzeugt die andere. Der
Zirkel löst sich auf. \enquote{Leer} hieße \enquote{keine Windung}, und die
Windung \emph{ist}, was Zeit ist.

\textbf{[S]} \textbf{Die Asymmetrie, ehrlich benannt.} Die beiden Ränder sind
nicht gleich gut bestimmt. Der \emph{dichte} Pol ist hart und abgeleitet: $\tau_0$
und $L_0=\xi\,\ell_P$ folgen aus $\xi$. Der \emph{dünne} Pol ist strukturell
ausgeschlossen, aber \emph{unbeziffert}: dass $m=0$ unmöglich ist, folgt aus der
Relation; \emph{wie klein} $m$ minimal wird, folgt daraus nicht -- und zwar aus
einem strukturellen Grund. Aus $T\cdot m=1$ trägt jedes System seine eigene
Obergrenze $R_\text{Torus}(m)=\hbar/mc$; eine universelle Maximalgröße ist nicht
sinnvoll gefordert. Für Systeme jenseits der Teilchenskala müsste daher eine
eigene, unabhängig gesetzte Obergrenze eingeführt werden -- das ist keine Lücke,
sondern Folge der Relation, und es liegt außerhalb dieses Dokuments (A010).

\section{Die Randfrage}

\textbf{[B]} Der Torus ist kompakt und randlos, $\partial T^4=\emptyset$ (A090).
Die Frage \enquote{was war bei $t=0$?} setzt einen Rand voraus, an dem ein Anfang
sitzen könnte; ein kompakter Torus hat keinen. Die Frage ist damit nicht falsch
\emph{beantwortet}, sondern falsch \emph{gestellt}: sie importiert die
dekompaktifizierte Projektion als Ontologie. Was eine ausgerollte Beschreibung als
singulären Anfang liest, ist der Punkt, an dem die \emph{Projektion} endet -- nicht
die Welt.

\section{Die Projektions-Einsicht}

\textbf{[S]} Es liegt nahe, die Systeme auf einer gemeinsamen Achse anzuordnen --
von maximaler zu minimaler Massenkonzentration -- und sich selbst in deren Mitte
zu finden. Diese Achse ist jedoch eine \emph{Außenbetrachtung}: sie setzt einen
Standpunkt voraus, von dem aus alle Systeme nebeneinander erscheinen. Einen
solchen Standpunkt gibt es in FFGFT nicht. $T\cdot m=1$ gilt \emph{je System}; es
gibt kein gemeinsames $T$, gegen das aufgetragen werden könnte. Schärfer: wir
\emph{projizieren uns in die Mitte} -- die Achse entsteht mit dem Betrachter als
Nullpunkt, und die Mitte ist keine Entdeckung, sondern eine Setzung, die
unbemerkt bleibt, weil der Setzende und der Verortete dasselbe sind. Die Ränder
erscheinen als \enquote{extrem}, weil sie vom gewählten Nullpunkt weit entfernt
liegen: \emph{extrem} heißt \emph{fern}, nicht \emph{ausgezeichnet}.

\textbf{[S]} Die intrinsische Zeit bleibt davon unberührt. Für ein Elektron ist
ohne Belang, wo ein Außenstehender es einträgt: seine Windung ist dieselbe, denn
$T=1/m$ und $m$ ist dasselbe. Was sich unterscheidet, ist allein, wie ein anderes
System es sieht -- eine Relation zwischen zweien, keine Eigenschaft des einen.
Auch die \enquote{langsamste Zeit} am Horizont ist eine Aussage über den
Vergleich, nicht über das Erleben. Die Projektion ist damit nicht ungültig,
sondern \emph{Werkzeug}: ohne gemeinsame Achse ließe sich nichts vergleichen. Sie
ist dieselbe Bewegung, die schon die Frage nach $t=0$ auflöste -- diesmal auf die
Skala selbst angewandt.

\section{Der Grundtakt}

\textbf{[K]} Die zulässigen Zustände sind diskrete Moden, und sie stehen
untereinander stets in fester Relation -- ihre Verhältnisse sind exakt. Daraus
folgt ein gemeinsamer \emph{Grundtakt}: über $\tilde T\cdot m=1$ trägt zwar jede
massive Größe ihren eigenen Zyklus $T=\hbar/mc^2$, doch diese Zyklen sind keine
unabhängigen Uhren, sondern kommensurable Vielfache desselben $\xi$-fixierten
Inkrements $d_1=1/75$ -- \emph{Harmonische eines Grundtakts}, nicht ein Kontinuum
freier Uhren. Das ist die präzise Fassung: nicht eine ausgezeichnete Grundeinheit,
gegen die alles gemessen würde, sondern ein gemeinsamer Takt, dessen Vielfache die
Moden sind.

\section{Der Zeitpfeil}

\textbf{[S]} Warum vergeht die Zeit in eine Richtung? In FFGFT ist die
Zeitrichtung \emph{geometrisch}: sie folgt aus der Orientierung der Windung auf
der kompakten Zeit-Richtung, nicht aus der Thermodynamik. Der thermodynamische
Zeitpfeil ist \emph{sekundär} -- eine emergente Eigenschaft der nicht-kompakten
Raumanteile, keine fundamentale Vorgabe. Das erklärt die scheinbare
Irreversibilität der makroskopischen Physik, ohne auf Ad-hoc-Annahmen über
Anfangsbedingungen zurückzugreifen. Und die Zeitrichtung ist \emph{skalenabhängig}:
auf der fundamentalen Ebene ist sie geometrisch fixiert, auf der makroskopischen
erscheint sie als statistischer Gradient. Beide Ebenen widersprechen sich nicht --
die eine ist die Projektion der anderen.

\section{Granulation und fundamentale Asymmetrie}

\textbf{[S]} Der tiefere Grund, dass überhaupt Zeit \enquote{vergeht}, ist eine
\emph{fundamentale Asymmetrie} als Bewegungsprinzip: die Windung ist nicht
symmetrisch, und die Abweichung -- der Defekt $100\,\xi_n$ pro Schritt -- ist der
Motor, der die Zeit erzeugt. Die Realität ist dabei \emph{granular}: unterhalb der
fundamentalen Skala gibt es kein Kontinuum, sondern diskrete Inkremente; erst
oberhalb bestimmter Skalen \emph{sieht} die Zeit stetig aus, weil die Granulation
$d_1=1/75$ relativ zu makroskopischen Zeiten verschwindet. Das ist der Grund,
warum die praktische Physik mit kontinuierlicher Zeit korrekt rechnet, obwohl die
fundamentale Struktur diskret ist: die Stetigkeit ist eine gute Näherung, keine
Grundeigenschaft.

\section{Wer „Zeit nicht extern`` bereits bemüht hat}

\textbf{[S]} Der Gedanke, die Zeit nicht als äußeren Parameter zu behandeln, ist
keine Neuheit, sondern eine lange verfolgte Linie -- und es ist redlich, die
eingerollte Sicht in ihr zu verorten und dagegen abzugrenzen. \textbf{Relationale
Zeit:} Page \& Wootters lassen die Zeit \emph{innerhalb} des Zustandsraums
entstehen, als Verschränkung zwischen Uhr-Subsystem und Rest (\enquote{Evolution
ohne Evolution}); Moreva et al.\ und Marletto \& Vedral haben das präzisiert und
illustriert; Gambini \& Pullin entwickelten die Montevideo-Interpretation.
\textbf{Zeitlose Quantengravitation:} Wheeler--DeWitt ($\hat H\Psi=0$, das
\enquote{Problem der Zeit}); Kuchař und Isham lieferten die Übersichten; Rovelli
(relationale QM, \enquote{Forget Time}, mit Connes die thermische-Zeit-Hypothese);
Barbour (\emph{The End of Time}). \textbf{Kompakte Zeit:} in der Stringtheorie
untersucht; Bars entwickelte eine \enquote{zweizeitige Physik}. \textbf{Moden als
Struktur-Eigenwert:} Wigner (Masse als Casimir-Invariante, ein Label
irreduzibler Poincaré-Darstellungen); Kaluza/Klein (eine kompaktifizierte
Dimension erzeugt einen diskreten Moden-Turm mit gequantelten Ladungen).

\textbf{[S]} Die Abgrenzung ist präzise: der überwiegende Teil dieser Tradition
macht die Zeit \emph{relational}, \emph{emergent} oder \emph{zeitlos} -- sie nimmt
die Zeit als Fundamentales \emph{heraus}. Die eingerollte Sicht tut etwas subtil
anderes: sie nimmt die Zeit \emph{hinein} -- als kompakte geometrische Dimension,
die Moden trägt. Die Neuheit liegt nicht in \enquote{Zeit ist nicht extern} (das
ist alt und gut belegt), sondern in der spezifischen Route \enquote{kompakte Zeit
erzeugt Massen-Moden}.

\section{Die absolute Skala: fixiert, nicht frei}

\textbf{[B]} Die Verhältnisse der Massen folgen aus der Modenstruktur und sind
exakt; die absolute Skala ist damit noch nicht in SI-Zahlen ausgedrückt -- aber
sie ist deshalb kein \emph{freier} Parameter. Das Verhältnisgerüst ruht auf dem
einen dimensionslosen $\xi=4/30000$. Die Absolutwerte tragen ihre Korrektur nicht
als Fit, sondern über feste \emph{Brückenkonstanten}: $v$ in $m=r\,\xi^{p}v$
(A100), $E_0$ in $\alpha=\xi E_0^2$ (A130), feste Faktoren in $G$. In Verhältnissen
kürzen diese weg ($m_\mu/m_e=206{,}768$, korrekturfrei). $\hbar$ und $c$ sind reine
SI-Umrechnung; ein \emph{gemessener} dimensionaler Anker im Standardsinn existiert
nicht.

\textbf{[K]} Die Planck-Größen sind kein fundamentaler Eingang, sondern folgen aus
$\xi$: die Kette lautet $\xi\to\{m_e,\dots\}\to G\to\ell_P$. Die \emph{einzige}
exakte $\xi$-Potenz-Relation ist dabei
\[
  \frac{L_0}{\ell_P} = \xi,\qquad L_0 = \xi\,\ell_P \approx 2{,}15\times10^{-39}\,\text{m};
\]
der SI-Faktor $\ell_P$ darf \emph{nicht} in einen $\xi$-Exponenten hineingerechnet
werden. Die Kausalrichtung ist also umgekehrt: die Planck-Größen folgen aus $\xi$,
nicht $\xi$ aus ihnen. Über $\tilde T\cdot m=1$ liegt der Zyklus fest, sobald die
Masse über die $\xi$-Geometrie feststeht -- es bleibt kein zu eichender Modul
offen.

\section{Kausalität: kompakte Zeit ist spektral, nicht kausal geschlossen}

\textbf{[B]} Die naheliegende Sorge -- erzeugt eine periodische Zeit geschlossene
zeitartige Kurven? -- greift nicht. Der Träger ist ein Torus ($\pi_1\neq0$, stabile
Windungen), kein einfacher Kreis; die relevante Verbindung ist nicht ontologische
Nichtlokalität, sondern eine \emph{topologische} Verbindung auf $T^4$ (A070). Die
klassische CTC-Problematik (Gödel-Lösungen, Chronologie-Schutz) betrifft den
\emph{dekompaktifizierten} Lorentz-Sektor; die kompakte Masse-Zeit-Richtung ist
\emph{spektral} geschlossen, nicht kausal.

\section{Prüfskripte}

\begin{itemize}[leftmargin=2em,itemsep=2pt,topsep=2pt]
  \item Native Reziprozität und Extrema: numerische Gegenprobe von
    $T\cdot E=1$, der exakten Schließung $\tau_0\,E_\text{max}=t_P E_P=1$ und
    $d_1=100\xi_1=1/75$:
    \texttt{python/A\_Serie\_Skripte/a060\_reziprozitaet.py}
  \item Grundtakt: dass die Moden-Zyklen kommensurable Vielfache von $d_1$ sind
    (ganzzahlige Verhältnisse), nicht freie Uhren:
    \texttt{python/A\_Serie\_Skripte/a060\_grundtakt.py}
\end{itemize}

\section{Was offen bleibt}

\textbf{Erstens ist die Länge der kompakten Zeit-Richtung nicht als eine SI-Zahl
gegeben.} Dass die Verhältnisse exakt sind und die absolute Skala über
Brückenkonstanten fixiert (nicht frei) ist, ist der Stand; ein gemessener
dimensionaler Anker im Standardsinn existiert nicht.

\textbf{Zweitens ist die \emph{quantum equivalence} der ausgerollten und der
eingerollten Sicht bekannt, ihre methodologischen Konsequenzen aber nur
systematisiert, nicht abschließend bewertet.} Dass beide denselben Formalismus
benutzen, ist gesichert; welche Lesart vorzuziehen ist, bleibt eine Setzung mit
ontologischem Preis.

\textbf{Drittens ist der dünne Pol strukturell ausgeschlossen, aber unbeziffert.}
Für Systeme jenseits der Teilchenskala verlangt die Relation eine eigene,
unabhängig gesetzte Obergrenze; deren Ausführung liegt außerhalb dieses Dokuments.

\section{Ersetzt}

Dieses Dokument nimmt auf: Dok.~010 (Einheitenkonvention, Zeit-Energie-Dualität,
universelle Feldgleichung, Lagrange-Dichte, charakteristische Längen und Energie,
Skalenhierarchie), Dok.~069 (Zeit-Energie-Dualität, geometrische
Ruhemasse-Definition), Dok.~078 (Granulation, fundamentale Asymmetrie,
Zeitstruktur), Dok.~157 (Zeitpfeil), Dok.~296 (Zeit, Offenheit, Invarianz;
Beobachter- und Projektionsteil), Dok.~306 (native Zeit-Energie-Reziprozität) und
Dok.~307 (Zeit im Zustandsraum, Grundtakt). Eingearbeitet: R50 (die
Heisenberg-basierte Singularitäts-Argumentation aus Dok.~025/063 ist durch die
native Kette $T\cdot E=1$ ersetzt -- Konklusion unverändert, Ableitung nativ) und
R54 (kein $N_0$ als Grundeinheit; die Moden sind kommensurable Vielfache eines
gemeinsamen Grundtakts). \emph{Nicht} übernommen: der kosmisch-große Sektor (der
statische $\xi$-Kosmos und die Hintergrundstrahlungs-Abschnitte aus Dok.~069, die
spekulativen Abschnitte aus
Dok.~296 §5, R46) sowie die Skalen-Schreibweise $E_H=E_0\xi^{41/4}$ (an eine
extern gemessene Größe kalibriert, keine Ableitung -- P35); dieser Bereich wird in der A-Serie nicht
geführt (A010) und verlangt eine eigene Skala. Der alte $T\cdot m=1$-Singularitäts\-beweis
aus Dok.~078 ist durch die native Kette (R50) abgelöst.

\section{Verwendung von KI-Werkzeugen}

Dieses Dokument ist unter Verwendung KI-gestützter Werkzeuge entstanden. Die
Fragestellungen, die Setzungen und alle inhaltlichen Entscheidungen stammen vom
Autor; die Werkzeuge formulieren aus, rechnen nach und erstellen Prüfskripte.
Die Verantwortung für Inhalt und Fehler liegt vollständig beim Autor. Der
ausführliche Wortlaut steht in A010.
